\documentclass[11pt]{article}
\usepackage[T1]{fontenc}
\usepackage[utf8]{inputenc}
\usepackage[french]{babel}
\usepackage[a4paper,margin=1in]{geometry}
\usepackage{lmodern}
\usepackage{microtype}
\usepackage{amsmath,amssymb,mathtools}
\usepackage{enumitem}
\setlength{\parindent}{0pt}
\setlength{\parskip}{0.55em}
\newcommand{\bbC}{\mathbb{C}}
\newcommand{\End}{\operatorname{End}}
\newcommand{\Lie}{\operatorname{Lie}}
\newcommand{\Sp}{\operatorname{Sp}}
\newcommand{\Ker}{\operatorname{Ker}}
\newcommand{\Imop}{\operatorname{Im}}
\newcommand{\ad}{\operatorname{ad}}
\newcommand{\expop}{\operatorname{exp}}
\begin{document}

\begin{flushleft}
1987
\end{flushleft}

\begin{center}
\textbf{Lettre à Millson}
\end{center}

Cher Millson,

Nori m'a apporté une copie de votre article avec Goldman sur les déformations de fibrés plats. Il me semble qu'il contient un certain nombre d'inexactitudes, mais aussi qu'une fois corrigé, il donne les résultats plus forts que vous demandez p. 8, lignes -10 à -7. Je commence par ceux-ci.

La philosophie, dont je ne m'étais pas rendu compte avant de lire votre article, semble être la suivante : en caractéristique zéro, un problème de déformation est contrôlé par une algèbre de Lie différentielle graduée, et des algèbres de Lie DG quasi-isomorphes donnent la même théorie de déformation. Si l'algèbre de Lie DG qui contrôle un problème est « formelle », c'est-à-dire quasi-isomorphe à $\bigoplus H^i$, alors l'espace de déformation formelle verselle est celui de $\bigoplus H^i$, c'est-à-dire le complété en $0$ du sous-schéma de $H^1$ défini par les équations $[u,u]=0$.

Si l'on croit cela, ce que vous faites devient transparent.

\paragraph{Problème de déformation.}
$G$ est un groupe réductif donné, et l'on veut déformer un $G$-torseur $P$. Plus précisément, on veut déformer un $G(\bbC)$-torseur, ou, de façon équivalente, un $G(\mathcal O)$-torseur, muni d'une connexion intégrable $\nabla$.

\paragraph{Algèbre DG contrôlante.}
Sur $\bbC$, l'algèbre DG contrôlante est
\[
  \Omega^{\bullet,\bullet}\bigl(\Lie(G^P)\bigr).
\]
Si $P$ peut être réduit à un sous-groupe compact maximal $K\subset G$, c'est-à-dire à une forme compacte de $G$, la théorie de Hodge montre que l'algèbre contrôlante est formelle : on a des quasi-isomorphismes d'algèbres de Lie DG
\begin{equation}\label{eq:quasi}
  \Omega^{\bullet,\bullet}\bigl(\Lie(G^P)\bigr)
  \longleftarrow
  \Ker(d')
  \longrightarrow
  \Ker(d')/\Imop(d'),
\end{equation}
et $\Ker(d')/\Imop(d')$ a $d=0$.

Pour vérifier ce qui précède, il suffit d'utiliser le principe des deux types, sous la forme suivante : le complexe double $\Omega^{\bullet,\bullet}(\Lie(G^P))$, en oubliant le crochet, est une somme de
\begin{enumerate}[label=\alph*)]
\item un complexe double avec $d'=d''=0$ (la « partie harmonique ») ;
\item carrés d'isomorphismes
\[
\begin{array}{ccc}
 * & \xrightarrow{\ d'\ } & * \\
 {\scriptstyle d''}\uparrow && \uparrow{\scriptstyle d''} \\
 * & \xrightarrow{\ d'\ } & *
\end{array}
\qquad (0\text{ ailleurs}).
\]
\end{enumerate}

En ce qui concerne les déformations formelles, cela termine l'histoire. Il ne m'est pas clair comment utiliser la théorie de Hodge dans le cas analytique pour obtenir le même résultat pour les espaces analytico-versels. Le cas algébrique va bien, mais Artin est plus puissant que cela.

Voici comment donner un sens à cette philosophie, et la vérifier.

\paragraph{Construction.}
Associer à une algèbre de Lie DG $L^*$ sur un corps $k$ de caractéristique zéro une catégorie fibrée au-dessus de la catégorie des spectres $\Sp(A)$, où $A$ est une $k$-algèbre locale artinienne de dimension finie, de corps résiduel $k$.

Les objets au-dessus de $\Sp(A)$, avec $k=A/\mathfrak m$, sont
\[
  \omega\in L^1\otimes\mathfrak m,
  \qquad
  d\omega+\frac12[\omega,\omega]=0.
\]

\paragraph{Flèches.}
L'algèbre de Lie $L^0$ donne naissance à un groupe formel $\widehat L^0$ ; les points à valeurs dans $A$ de ce groupe formel sont $L^0\otimes\mathfrak m$, vu comme algèbre de Lie nilpotente et muni de la loi de groupe de Campbell-Hausdorff. On note ce groupe $\widehat L^0(A)$.

Le groupe formel $\widehat L^0$ agit sur $L^1$ par le crochet $[\ ,\ ]:L^0\otimes L^1\to L^1$, donc $\widehat L^0(A)$ agit sur $L^1\otimes A$. On dispose aussi d'une application
\[
  g^{-1}dg:\widehat L^0\longrightarrow L^1,
  \qquad
  \expop(\lambda)\longmapsto
  \frac{1-\expop(-\ad\lambda)}{\ad\lambda}\,d\lambda,
\]
et donc d'une application $\widehat L^0(A)\to L^1\otimes\mathfrak m$.

On définit alors une application de $\omega\in L^1\otimes\mathfrak m$ vers $\omega'\in L^1\otimes\mathfrak m$ en prenant $g\in\widehat L^0(A)$ et en posant
\begin{equation}\label{eq:gauge}
  \omega'=g\,dg^{-1}+g(\omega).
\end{equation}
Il faut vérifier que
\[
  d\omega'+\frac12[\omega',\omega']
  =g\left(d\omega+\frac12[\omega,\omega]\right)
\]
lorsque $\omega$ et $\omega'$ sont reliés par \eqref{eq:gauge}.

\paragraph{À vérifier.}
Pour un idéal $I$ de carré nul dans $A$, et un objet $\omega$ sur $A/I$, il y a une obstruction dans $H^2(L\otimes I)$ à étendre $\omega$ à $A$. Si l'obstruction s'annule, les classes d'isomorphie d'extensions forment un espace homogène principal sous $H^1(L\otimes I)$, et les automorphismes de toute extension, triviaux sur $A/I$, sont $H^0(L^*\otimes I)$. Tout cela est fonctoriel en $L^*$, d'où le fait que les quasi-isomorphismes induisent des équivalences de catégories fibrées.

\paragraph{Application ici.}
Soit $X$ une variété analytique, et soit $(V,\nabla)$ un fibré vectoriel holomorphe muni d'une connexion plate. Le problème de déformation est le foncteur qui à $A$ associe la catégorie des fibrés vectoriels avec connexion dans la direction de $X$, sur $X\otimes\Sp(A)$, déformant $(V,\nabla)$.

Il est donné par la construction précédente appliquée à
\[
  \Omega^{\bullet,\bullet}(\End V).
\]

\paragraph{Remarque.}
Dans votre application, puisque vous disposez du quasi-isomorphisme explicite \eqref{eq:quasi}, vous avez une équivalence explicite de catégories fibrées entre
\begin{itemize}
\item les déformations de $(V,\nabla)$ comme ci-dessus ;
\item la catégorie dont les objets au-dessus de $\Sp(A)$ sont les applications pointées
\[
  \Sp(A)\longrightarrow \{u\in H^1_L\mid [u,u]=0\},
\]
et dont les flèches sont les éléments $g\in \widehat{H}^0_L(A)$ envoyant $u$ sur $g(u)$.
\end{itemize}

Le foncteur vers le haut est le suivant. Une classe de cohomologie $u\in H^1L\otimes\mathfrak m$, c'est-à-dire une application $u:\Sp(A)\to H^1L$, s'écrit de façon unique comme un $\omega\in L^1\otimes\mathfrak m$ avec $d'\omega=d''\omega=0$. Ce $\omega$ peut être corrigé par un $d'\rho$ pour avoir une courbure nulle, si $[u,u]=0$. Le $\widetilde\omega$ corrigé est unique modulo $\exp(\Ker d'\otimes\mathfrak m)$.

\paragraph{Remarque.}
Le quasi-isomorphisme \eqref{eq:quasi} est un quasi-isomorphisme filtré pour la filtration de Hodge. Cela devrait se traduire par l'assertion que le couple d'espaces
\[
  \bigl(\text{déformations de }(V,\nabla)\bigr)
  \supset
  \bigl(\text{déformations de }V\text{ sur }X\text{ fixé}\bigr)
\]
est isomorphe au couple
\[
  \bigl(H^1\supset F^1H^1\bigr)\cap\{[u,u]=0\}.
\]

\section*{II. Commentaires sur votre article}

p. 1, ligne 2 après « in B » : non. Prendre $G=U_1$.

p. 6, ligne 18, « Philosophic? » : je suppose que vous voulez commencer avec $\eta$ tel que
\[
  d'\eta=d''\eta=0.
\]

p. 3, ligne 4 : pourquoi « fermé » implique-t-il que la composante $(2,0)$ soit fermée ?

\section*{III. Généralisation}

La théorie de Hodge est disponible pour les variations de structure de Hodge, variations complexes, sur les variétés kählériennes compactes. Le bigradué utile de $\Omega^{\bullet,\bullet}(\End V)$ s'obtient en mélangeant celui de $\Omega$ et celui de $\End V$ ; il est utile maintenant que \eqref{eq:quasi} soit un quasi-isomorphisme filtré. On obtient :

\begin{enumerate}[label=\alph*)]
\item le même résultat qu'avant sur les déformations d'une représentation de $\pi_1$ donnée par une variation complexe ;
\item en regardant la théorie de déformation contrôlée par $F^0(\cdots)$, le même résultat pour le module versel d'un fibré filtré avec connexion provenant d'une variation complexe. On impose que la connexion déformée continue à satisfaire la transversalité :
\[
  \nabla F^i\subset \Omega^1\otimes F^{i-1}.
\]
\end{enumerate}

Bien sincèrement,

\begin{flushright}
P. Deligne
\end{flushright}

\end{document}
